\documentclass[12pt, letterpaper]{article}

%% Build from the paper/ directory:
%%   cd paper && pdflatex main.tex
%% or use the root build system:
%%   make paper          (from repo root)
%%   ./build.sh paper    (POSIX)
%%   build.cmd paper     (Windows cmd)

\input{macros/packages}
\input{macros/commands}

%% Page geometry (also set here as a fallback in case the package path fails)
\geometry{top=1in, bottom=1in, left=1in, right=1in}

%% ---------------------------------------------------------------
%% TITLE PAGE
%% ---------------------------------------------------------------
%% TODO: replace title, subtitle, author, ORCID, email.
%% TODO: at release time, fill in the \thanks{} block with version,
%%       Zenodo DOI, GitHub URL, and any series identifier.

\title{The Electric Induced-Action Block:\\
       Permittivity and the Covariant Completion of the Lattice Gauge Response%
       \thanks{Paper VIII of the A=1 Discrete Causal Lattice series.
       Version 0.1-DRAFT.
       %% Fill in at release time:
       %% \href{https://doi.org/10.5281/zenodo.XXXXXXXX}%
       %%      {DOI: 10.5281/zenodo.XXXXXXXX}.
       %% \href{https://github.com/JackDMenendez/dcl-paper-08-electric-induced-action}%
       %%      {GitHub: JackDMenendez/dcl-paper-08-electric-induced-action}.
       Pre-release draft; DOI to be assigned.}}
\author{Jack D. Menendez\\
        \small Eugene, OR\\
        \small ORCID: \href{https://orcid.org/0009-0003-1166-307X}{0009-0003-1166-307X}\\
        \small \href{mailto:jackdmenendez@gmail.com}{jackdmenendez@gmail.com}}
\date{\today}

\begin{document}

\maketitle

%% ---------------------------------------------------------------
%% FRONT-MATTER NOTICES
%% ---------------------------------------------------------------
\noindent\textbf{License.}  Paper text and figures: Creative Commons
Attribution 4.0 (CC BY 4.0).  Source code in the companion
repository: MIT License.

\medskip
\noindent\textbf{Data availability.}  Numerical experiments, data
files, and the LaTeX source for this paper are publicly available at
\href{https://github.com/JackDMenendez/dcl-paper-08-electric-induced-action}%
     {\nolinkurl{github.com/JackDMenendez/dcl-paper-08-electric-induced-action}}.
%% At release time, point to the tag and the Zenodo DOI:
%% The exact commit corresponding to this release is the vX.Y tag,
%% archived at
%% \href{https://doi.org/10.5281/zenodo.XXXXXXXX}%
%%      {\nolinkurl{doi:10.5281/zenodo.XXXXXXXX}}.

\medskip
\noindent\textbf{Keywords.}  discrete causal lattice; induced gauge action;
permittivity; vacuum birefringence; Peierls substitution; Wilson loop; A=1.

\newpage

\begin{abstract}
    % abstract.tex -- Paper VIII
On the A=1 Discrete Causal Lattice the magnetic field enters the tick rule as a
spatial link phase and yields, through the bipartite-plaquette Wilson loop, an exact
induced-action tensor with eigenvalues $\{4,4,16\}$ and optical axis $(1,1,-1)$
(Paper~I, Appendix~B). The electric field enters differently---as an on-site,
mass-like phase advance---so it admits no Wilson-loop analogue, and the corresponding
electric action block (the permittivity $\varepsilon$) is absent from both Paper~I and
Paper~II. Because the framework establishes only the spatial point group $O_h$ and not
the Lorentz boosts, there is no symmetry route from the magnetic $\mu^{-1}$ to the
electric $\varepsilon$: the electric block must be derived directly. We carry out that
derivation and assemble the covariant $(\varepsilon,\mu^{-1})$ completion of the
lattice gauge response. The result decides a binary physical question left open by
Paper~IV: whether the lattice's vacuum birefringence about $(1,1,-1)$ cancels---so the
framework predicts the \emph{absence} of gauge-sector birefringence---or persists as a
genuine, testable signal. The magnetic $\{4,4,16\}$ result is retained throughout as an
exact consistency anchor.

\end{abstract}

%% ---------------------------------------------------------------
%% RELEASE NOTES (only at release time)
%% ---------------------------------------------------------------
%% \bigskip
%% \noindent\textbf{Release notes for vX.Y (since vX.Y-1).}
%% \begin{itemize}
%% \item Item 1.
%% \item Item 2.
%% \end{itemize}

%% \emph{Status legend used throughout.}  \texttt{PASS} / \texttt{PART}
%% / \texttt{STUB} / \texttt{FAIL} -- see Table~\ref{tab:audit}.

\newpage
\tableofcontents
\newpage

%% ============================================================
%% PART I: <name>
%% ============================================================
\section{Introduction}
\input{sections/introduction}

\section{A Section Showing the Patterns}
\input{sections/section_template}

%% Add more \section{...} \input{sections/...} blocks as needed.
%% Group with comment dividers (e.g., "PART II: DYNAMICS") to keep
%% main.tex skimmable as the paper grows.

\section{Conclusion}
\input{sections/conclusion}

\section*{Acknowledgements}
\input{sections/acknowledgements}

%% ============================================================
%% APPENDICES
%% ============================================================
\appendix

\section{Code and Data}
\input{sections/code_and_data}

\section{Reproducibility Quickstart}
\input{sections/reproducibility}

%% ============================================================
\bibliographystyle{unsrt}
\bibliography{paper-bib/references}

\end{document}
